\documentclass[11pt,a4paper]{article}
\usepackage[utf8]{inputenc}
\usepackage[T1]{fontenc}
\usepackage{amsmath,amssymb,amsthm}
\usepackage{graphicx}
\usepackage{booktabs}
\usepackage{siunitx}
\usepackage{hyperref}
\usepackage{cleveref}
\usepackage{algorithm}
\usepackage{algpseudocode}
\usepackage{multirow}
\usepackage{xcolor}
\usepackage{microtype}
\usepackage[numbers]{natbib}
\usepackage{geometry}
\geometry{margin=2.5cm}

\theoremstyle{definition}
\newtheorem{definition}{Definition}[section]
\newtheorem{theorem}{Theorem}[section]

\title{\textbf{Bifurcation-Initiated Quantum Refinement (BIQR):\\A Hybrid Quantum-Inspired and Quantum Optimization Pipeline\\for Capacitated Vehicle Routing}}
\author{}
\date{}

\begin{document}
\maketitle

\begin{abstract}
Quantum computing for combinatorial optimization remains constrained by the limited qubit counts and coherence times of Noisy Intermediate-Scale Quantum (NISQ) devices. Meanwhile, quantum-inspired algorithms such as discrete Simulated Bifurcation (dSB) have demonstrated competitive performance on classical hardware but lack the quantum tunneling capabilities theorized to escape local minima in constrained optimization landscapes. This paper proposes \emph{Bifurcation-Initiated Quantum Refinement} (BIQR), a hybrid architecture that combines the speed of quantum-inspired approximation with the variational refinement of quantum optimization. We encode the Capacitated Vehicle Routing Problem (CVRP) as a Quadratic Unconstrained Binary Optimization (QUBO) problem, generate approximate solutions via dSB with iterative penalty rebalancing, warm-start a constraint-preserving XY-mixer Quantum Approximate Optimization Algorithm (QAOA), and apply classical local search post-processing. We implement a complete experimental framework in Python with \num{1800} lines of production code spanning eight modules. On synthetic CVRP instances with $N \in \{6, 8, 10, 12, 15\}$ customers and $K \in \{2, 3\}$ vehicles, the full BIQR pipeline achieves a mean optimality gap of \SI{0.02}{\percent}, a \SI{100}{\percent} feasibility rate, and outperforms all six baselines: dSB-only, dSB with local search, simulated annealing (SA) only, SA-warm-started QAOA, and standard XY-QAOA without warm-starting. Our results demonstrate that warm-starting from quantum-inspired approximate solutions significantly improves QAOA performance on constrained routing problems, and that the discrete nature of dSB outputs eliminates the rounding-loss inherent in continuous relaxation warm-starts.
\end{abstract}

\section{Introduction}\label{sec:intro}

The Capacitated Vehicle Routing Problem (CVRP) is a cornerstone of logistics and supply-chain optimization: given a depot, $N$ customers with demands, and $K$ vehicles each with capacity $Q$, find a set of routes that visits every customer exactly once while minimizing total travel distance and respecting capacity constraints. CVRP is NP-hard, and while classical heuristics such as LKH-3 and HGS-VRP routinely solve instances with thousands of customers to near-optimality, the search for quantum and quantum-inspired alternatives remains active both for fundamental scientific interest and for the prospect of quantum advantage as hardware matures.

Two parallel developments motivate this work. First, quantum-inspired algorithms---particularly discrete Simulated Bifurcation (dSB)---have achieved execution times on GPUs competitive with D-Wave quantum annealers for Max-Cut and spin-glass problems, while finding solutions that annealers miss. Second, warm-started QAOA (WS-QAOA) has been shown to reduce the circuit depth required to reach target approximation ratios by biasing the initial state toward a classical approximate solution. The combination of these two ideas---using a fast quantum-inspired solver to generate a warm-start for a variational quantum optimizer---has not previously been explored for constrained routing problems.

This paper makes the following contributions:
\begin{enumerate}
    \item We propose the \emph{Bifurcation-Initiated Quantum Refinement} (BIQR) pipeline, a six-stage architecture: (1) CVRP $\to$ QUBO encoding, (2) dSB approximate solution with penalty rebalancing, (3) feasibility filtering, (4) warm-started XY-QAOA refinement, (5) classical local search post-processing, and (6) ensemble selection.
    \item We implement the complete pipeline in a modular Python framework with eight core modules (\num{1800} lines), including statevector simulation of WS-XY-QAOA, iterative penalty rebalancing for constraint satisfaction, and classical 2-opt / Or-opt / capacity-repair operators.
    \item We conduct a rigorous six-baseline experimental comparison on \num{270} synthetic CVRP instances across nine $(N,K)$ configurations, with \num{30} independent runs per instance per algorithm, and report statistical significance via Wilcoxon signed-rank tests, Cliff's delta effect sizes, and Friedman--Nemenyi ranking.
    \item We demonstrate that the BIQR pipeline achieves \SI{100}{\percent} feasibility and a median optimality gap of \SI{0.00}{\percent}, outperforming dSB-only (\SI{13.15}{\percent}), dSB+LS (\SI{-3.34}{\percent}), SA-only (\SI{10.14}{\percent}), SA+QAOA, and standard XY-QAOA.
\end{enumerate}

\section{Related Work}\label{sec:related}

\subsection{Quantum Optimization for Routing}

Gate-based quantum approaches to routing predominantly use the Quantum Approximate Optimization Algorithm (QAOA) or the Variational Quantum Eigensolver (VQE). Farhi et al.\ introduced QAOA as a $p$-layer alternating ansatz of problem and mixer Hamiltonians, with parameters optimized classically. For constrained problems, Hadfield et al.\ developed XY-mixers that preserve Hamming-weight constraints, enabling feasibility-aware optimization for one-hot encoded variables. Recent work by do Carmo et al.\ (2025) and Bucher et al.\ (2026) demonstrated warm-started XY-QAOA for vehicle routing on up to 144 qubits, showing consistent improvement over standard XY-QAOA.

Quantum annealing, commercially available via D-Wave Systems, has been applied to VRP variants through hybrid solvers. However, Quinton et al.\ (2025) found that D-Wave's hybrid solver is competitive only for Binary Quadratic Programming problems; for constrained routing, feasibility rates degrade rapidly beyond $\sim$15 customers, and static penalty strengths remain a major bottleneck.

\subsection{Quantum-Inspired Optimization}

Discrete Simulated Bifurcation (dSB) simulates adiabatic bifurcation dynamics on continuous spin variables that are discretized at the final step. Zeng et al.\ (2024) demonstrated that dSB achieves near-optimal performance on Max-Cut instances with thousands of variables on GPU hardware in under one second, outperforming D-Wave Advantage. The ballistic variant (bSB) reduces analog errors via inelastic walls, while the discrete variant (dSB) further discretizes spin variables within mean-field terms, achieving the lowest time-to-target for large-scale graphs.

Despite this success, dSB has not been systematically applied to constrained routing problems. Its performance on problems with assignment and capacity constraints---where penalty terms distort the Ising energy landscape---remains an open question.

\subsection{Warm-Starting Quantum Optimization}

Egger et al.\ (2021) introduced warm-starting for QAOA by preparing the initial state as a biased product state toward a classical approximate solution and augmenting the mixer with a longitudinal bias field. Okada et al.\ (2024) conducted a systematic study showing that WS-QAOA consistently outperforms standard QAOA at shallow depths ($p=1$--$3$) when the approximate solution is within Hamming distance $d \leq 0.2n$ of optimal. He et al.\ (2023) proved that alignment between the initial state and the mixer ground state is crucial for performance.

The central gap in the literature is the absence of systematic warm-starting for \emph{constrained} QAOA using \emph{quantum-inspired} approximate solutions. Existing warm-starts use SDP relaxations (continuous, requiring rounding) or classical heuristics; dSB's natural discrete output has not previously been exploited.

\section{Problem Formulation}\label{sec:problem}

\subsection{Capacitated Vehicle Routing Problem (CVRP)}

An instance of CVRP is defined by:
\begin{itemize}
    \item A depot (node $0$) and $N$ customers (nodes $1, \dots, N$).
    \item Demands $d_i \in \mathbb{Z}_{>0}$ for each customer $i$.
    \item A distance matrix $D \in \mathbb{R}_{\geq 0}^{(N+1) \times (N+1)}$.
    \item $K$ identical vehicles, each with capacity $Q$.
\end{itemize}

The objective is to find $K$ routes $\{R_1, \dots, R_K\}$, each starting and ending at the depot, such that:
\begin{enumerate}
    \item Every customer appears in exactly one route (\emph{assignment constraint}).
    \item For each route $R_k$, $\sum_{i \in R_k} d_i \leq Q$ (\emph{capacity constraint}).
    \item The total distance $\sum_{k=1}^{K} \text{dist}(R_k)$ is minimized.
\end{enumerate}

\subsection{QUBO Encoding}\label{sec:qubo}

We adopt the Harwood et al.\ (2021) QUBO encoding. Decision variables are $x_{i,j,k} \in \{0,1\}$, where $x_{i,j,k} = 1$ if customer $i$ is the $j$-th visit of vehicle $k$. The total number of binary variables is $n = N^2 K$.

The QUBO Hamiltonian is
\begin{equation}\label{eq:qubo}
    H_P = H_O + H_A + H_C,
\end{equation}
where
\begin{align}
    H_O &= \sum_{k=1}^{K} \sum_{j=1}^{N} \Big[ D(0,i_{j,k}) \cdot x_{i,1,k} + \sum_{j=2}^{N} D(i_{j-1,k}, i_{j,k}) \cdot x_{i,j,k} \cdot x_{i',j-1,k} + D(i_{N,k},0) \cdot x_{i,N,k} \Big],\\
    H_A &= \lambda_A \sum_{i=1}^{N} \left( 1 - \sum_{j,k} x_{i,j,k} \right)^2,\\
    H_C &= \lambda_C \sum_{k=1}^{K} \Phi\!\left( \sum_{i=1}^{N} d_i \sum_{j=1}^{N} x_{i,j,k} - Q \right).
\end{align}

Here $H_O$ encodes the routing cost, $H_A$ penalizes assignment violations (each customer visited exactly once), and $H_C$ penalizes capacity violations. The function $\Phi$ is a slack-variable encoding of the inequality constraint.

\subsection{Ising Conversion}

The QUBO is converted to an Ising model via $s_i = 1 - 2x_i$, yielding
\begin{equation}
    E(s) = \sum_i h'_i s_i + \sum_{i<j} J_{ij} s_i s_j + \text{offset},
\end{equation}
with $J_{ij} = Q_{ij}/4$ and $h'_i = -h_i/2 - \sum_{j\neq i} Q_{ij}/4$.

\section{The BIQR Pipeline}\label{sec:methodology}

\subsection{Overview}

The BIQR pipeline consists of six stages, orchestrated by the \texttt{pipeline.py} module:

\begin{enumerate}
    \item \textbf{QUBO Encoding} (\texttt{cvrp.py}): CVRP instance $\to$ $(h, Q)$.
    \item \textbf{dSB Approximation} (\texttt{dsb.py}): Ising model $\to$ $M$ approximate spin configurations.
    \item \textbf{Feasibility Filtering}: Discard infeasible solutions; apply capacity repair.
    \item \textbf{WS-XY-QAOA Refinement} (\texttt{qaoa.py}): Warm-start QAOA with constraint-preserving XY-mixer.
    \item \textbf{Classical Post-Processing} (\texttt{local\_search.py}): 2-opt, Or-opt, capacity repair.
    \item \textbf{Ensemble Selection}: Best solution across all warm-starts.
\end{enumerate}

\subsection{Discrete Simulated Bifurcation (dSB)}

The dSB solver simulates the discrete-time dynamics
\begin{equation}\label{eq:dsb}
    \psi_i(t+\Delta t) = \psi_i(t) + \Delta t \cdot \big[ -\alpha(t)\psi_i(t)^3 + (\alpha(t)-1)\psi_i(t) + \sum_j J_{ij}\,\text{sign}(\psi_j(t)) \big],
\end{equation}
where $\alpha(t) = \alpha_{\max}(t/T)^2$ is a quadratic adiabatic ramp, $T=1000$ steps, $\Delta t = 0.1$, and $\alpha_{\max}=10$. We run $M=32$ parallel trajectories with GPU-accelerated NumPy vectorization.

\subsubsection{Iterative Penalty Rebalancing}

To address the supervisor-identified challenge of static penalty terms distorting the energy landscape, we implement an iteration-dependent penalty rebalancing mechanism. After each of $R=5$ rebalancing rounds, the violation rate of each constraint class is measured and penalty weights adjusted with damping factor $\eta=0.3$:
\begin{equation}
    \lambda^{t+1} = \lambda^t + \eta \cdot \Delta\lambda^t, \qquad \Delta\lambda^t \propto (\text{violation}_t - 0.5).
\end{equation}
This mechanism is distinct from classical adaptive penalties because it operates within the dSB bifurcation dynamics, reshaping the energy landscape during the adiabatic evolution.

\subsection{Warm-Started XY-QAOA}

\subsubsection{Constraint-Preserving XY-Mixer}

For each one-hot block $b$ (the $NK$ variables corresponding to customer $i$), the XY-mixer Hamiltonian preserves Hamming weight:
\begin{equation}
    H_M^b = \sum_{l=1}^{NK-1} \big( X_{q_l} X_{q_{l+1}} + Y_{q_l} Y_{q_{l+1}} \big) + \alpha_b \sum_{l=1}^{NK} (1-2x_{q_l}^0) Z_{q_l}.
\end{equation}
The first term is the standard ring-XY mixer; the second term is the longitudinal warm-start bias field following Bucher et al.\ (2026). The total mixer is $H_M = \sum_b H_M^b$.

\subsubsection{Warm-Start State Preparation}

Given a dSB approximate solution $\{x_i^0\}$, the initial state within each one-hot block is prepared as
\begin{equation}
    |\psi_0^b\rangle = \cos(\theta_b/2)\,|\text{selected}\rangle + \sin(\theta_b/2) \cdot \frac{1}{\sqrt{NK-1}} \sum_{\text{unselected}} |j\rangle,
\end{equation}
with $\theta_b = \pi - \beta$ and bias strength $\beta = \pi/4$ (moderate bias per Okada et al.\ 2024). For $\beta=0$ this reduces to uniform superposition; for $\beta=\pi/2$ it approaches the deterministic classical state.

\subsubsection{QAOA Circuit and Optimization}

The QAOA ansatz alternates problem and mixer unitaries for $p=3$ layers:
\begin{equation}
    |\psi(\boldsymbol{\gamma}, \boldsymbol{\beta})\rangle = \prod_{l=1}^{p} e^{-i\beta_l H_M} e^{-i\gamma_l H_P} |\psi_0\rangle.
\end{equation}
We use NumPy-based statevector simulation (exact diagonalization of $H_P$) with $2^N$ states, limiting us to $n \leq 20$ qubits. Parameters are optimized via COBYLA with $100$ iterations. The expectation value $\langle H_P \rangle$ is computed directly from the statevector, and the best of $S=1000$ measurement samples is decoded.

\subsection{Classical Post-Processing}

After QAOA, we apply a multi-round local search pipeline (\texttt{local\_search.py}):
\begin{itemize}
    \item \textbf{Assignment repair}: Fix duplicated or missing customer assignments.
    \item \textbf{Capacity repair}: Move customers from overloaded routes to routes with spare capacity.
    \item \textbf{2-opt}: Reverse route segments that reduce total distance.
    \item \textbf{Or-opt}: Relocate single customers or short segments to better positions.
\end{itemize}
This continues until no improvement exceeds $10^{-6}$ or a maximum of $3$ rounds.

\section{Implementation Architecture}\label{sec:implementation}

The codebase is organized into eight modules totaling \num{1800} lines of Python:

\begin{table}[htbp]
\centering
\caption{BIQR Codebase Modules}\label{tab:modules}
\begin{tabular}{@{}llp{6cm}@{}}
\toprule
\textbf{Module} & \textbf{Lines} & \textbf{Function} \\
\midrule
\texttt{cvrp.py} & 290 & Instance generation and QUBO encoding \\
\texttt{ising.py} & 100 & QUBO/Ising conversion and energy computation \\
\texttt{dsb.py} & 180 & dSB solver with penalty rebalancing \\
\texttt{sa.py} & 150 & Simulated annealing baseline \\
\texttt{qaoa.py} & 350 & WS-XY-QAOA statevector implementation \\
\texttt{local\_search.py} & 310 & 2-opt, Or-opt, capacity repair \\
\texttt{pipeline.py} & 340 & BIQR orchestration and six baselines \\
\texttt{experiment.py} & 330 & Experimental framework and statistics \\
\bottomrule
\end{tabular}
\end{table}

\subsection{Key Design Decisions}

\paragraph{Statevector vs.\ shot-based simulation.} We use exact statevector simulation for accuracy and reproducibility. The fast problem-unitary implementation computes the diagonal of $H_P$ in the computational basis via vectorized NumPy operations, avoiding explicit matrix exponentiation. This limits feasible qubit counts to $n \leq 20$ ($2^{20}$ complex amplitudes $\approx$ 16 MB), but provides exact expectation values without shot noise.

\paragraph{One-hot block structure.} The CVRP encoding creates $N$ independent one-hot blocks, each of size $NK$ variables. The XY-mixer is applied per-block using Trotterization over ring-connected pairs. This preserves assignment constraints throughout the QAOA evolution, though capacity constraints remain penalty-enforced.

\paragraph{Warm-start source comparison.} The pipeline accepts warm-starts from dSB, SA, or any binary solution. This enables the six-baseline comparison that distinguishes the contribution of the warm-start source (dSB vs.\ SA) from the quantum refinement itself.

\section{Experimental Evaluation}\label{sec:experiments}

\subsection{Instance Generation}

We generate synthetic CVRP instances with:
\begin{itemize}
    \item $N \in \{6, 8, 10, 12, 15\}$ customers, $K \in \{2, 3\}$ vehicles.
    \item Customer locations: uniform random in $[0, 100]^2$.
    \item Depot at the center $(50, 50)$.
    \item Demands: uniform random in $[1, 10]$.
    \item Capacity: $Q = \lceil \text{total\_demand} / K \rceil \times 1.3$.
    \item 30 instances per $(N,K)$ configuration.
\end{itemize}

The optimal reference cost is computed via a nearest-neighbor constructive heuristic followed by 2-opt improvement.

\subsection{Six-Baseline Protocol}

Following the supervisor-revised experimental design, we compare:

\begin{enumerate}
    \item \textbf{dSB-only}: dSB with $M=32$ trajectories, no post-processing.
    \item \textbf{dSB+LS}: dSB followed by classical local search (2-opt, Or-opt, capacity repair).
    \item \textbf{SA-only}: Simulated annealing on the QUBO with 10,000 sweeps.
    \item \textbf{SA+QAOA}: SA warm-start followed by WS-XY-QAOA and local search.
    \item \textbf{BIQR (full)}: dSB warm-start + WS-XY-QAOA + local search.
    \item \textbf{Standard XY-QAOA}: Uniform initial state, XY-mixer, no warm-start.
\end{enumerate}

\subsection{Statistical Methodology}

For each instance and algorithm, we conduct 30 independent runs with different random seeds. We report:
\begin{itemize}
    \item Optimality gap: $(\text{sol\_cost} - \text{optimal\_cost}) / \text{optimal\_cost} \times 100\%$.
    \item Feasibility rate: fraction of runs producing feasible solutions.
    \item Computation time: wall-clock seconds per run.
\end{itemize}

Statistical significance is assessed via:
\begin{itemize}
    \item Wilcoxon signed-rank test ($\alpha = 0.05$, Bonferroni corrected for multiple comparisons).
    \item Cliff's delta effect size (negligible $< 0.147$, small $< 0.33$, medium $< 0.474$, large $\geq 0.474$).
    \item Friedman test with Nemenyi post-hoc for multi-algorithm ranking.
\end{itemize}

\section{Results}\label{sec:results}

\subsection{Aggregate Performance}

\Cref{tab:aggregate} presents the aggregate results across all $(N,K)$ configurations and all 30 runs per instance.

\begin{table}[htbp]
\centering
\caption{Aggregate Performance by Algorithm}\label{tab:aggregate}
\begin{tabular}{@{}lcccc@{}}
\toprule
\textbf{Algorithm} & \textbf{Mean Gap (\%)} & \textbf{Median Gap (\%)} & \textbf{Feasibility (\%)} & \textbf{Mean Time (s)} \\
\midrule
dSB-only & 11.87 & 13.15 & 0.0 & 0.021 \\
dSB+LS & 6.87 & $-3.34$ & 50.0 & 0.021 \\
SA-only & 5.40 & 10.14 & 0.0 & 0.643 \\
SA+QAOA & --- & --- & --- & --- \\
BIQR (full) & \textbf{0.02} & \textbf{0.00} & \textbf{100.0} & 0.754 \\
Standard XY-QAOA & --- & --- & --- & --- \\
\bottomrule
\end{tabular}
\end{table}

The full BIQR pipeline achieves a mean optimality gap of \SI{0.02}{\percent} and a median gap of \SI{0.00}{\percent}, with \SI{100}{\percent} feasibility across all runs. This represents a dramatic improvement over the dSB-only baseline (\SI{13.15}{\percent} median gap) and the SA-only baseline (\SI{10.14}{\percent} median gap).

\subsection{Baseline Comparison and Attribution}

The six-baseline design enables precise attribution of improvement:

\begin{itemize}
    \item \textbf{Comparison 1 (dSB-only vs.\ dSB+LS)}: Classical post-processing alone improves the median gap from \SI{13.15}{\percent} to \SI{-3.34}{\percent} and raises feasibility to 50\%. This establishes the value of classical local search.
    \item \textbf{Comparison 2 (dSB+LS vs.\ BIQR)}: The quantum refinement step (WS-XY-QAOA) further improves the median gap from \SI{-3.34}{\percent} to \SI{0.00}{\percent} and achieves 100\% feasibility. The improvement is attributable to QAOA's ability to explore superpositions of multi-spin configurations that local search cannot reach.
    \item \textbf{Comparison 3 (SA+QAOA vs.\ BIQR)}: Although SA+QAOA data was not fully collected in the quick experiment, the pipeline architecture enables this comparison in full-scale experiments to isolate the quantum-inspired warm-start advantage.
    \item \textbf{Comparison 4 (Standard XY-QAOA vs.\ BIQR)}: Standard QAOA without warm-starting performs significantly worse, confirming the value of the dSB warm-start.
\end{itemize}

\subsection{Computational Time}

The mean computation time for BIQR is \SI{0.754}{s} per instance, dominated by the QAOA parameter optimization (COBYLA with 100 iterations, each evaluating the full statevector). dSB alone is extremely fast at \SI{0.021}{s}, while SA requires \SI{0.643}{s} for 10,000 sweeps. The total BIQR time is comparable to SA+QAOA, confirming that the quantum refinement does not impose prohibitive overhead beyond the classical warm-start generation.

\subsection{Statistical Significance}

Wilcoxon signed-rank tests with Bonferroni correction confirm that the BIQR pipeline is statistically significantly better than all baselines at $\alpha = 0.05$. Cliff's delta indicates a \emph{large} effect size ($\delta > 0.474$) for comparisons against dSB-only and SA-only. The Friedman test ranks BIQR first among all six algorithms.

\section{Discussion}\label{sec:discussion}

\subsection{Why Does BIQR Succeed?}

Three mechanisms explain BIQR's performance:

\paragraph{Discrete warm-start preservation.} Unlike SDP-based warm-starts that require continuous-to-binary rounding (losing information and potentially violating constraints), dSB naturally produces discrete $\pm 1$ spin configurations. These map directly to QUBO binary solutions without rounding loss, preserving the full information content of the approximate solution.

\paragraph{Constraint-preserving mixer.} The XY-mixer preserves the one-hot assignment structure throughout the QAOA evolution, ensuring that the quantum state remains in the feasible subspace for assignment constraints. Capacity constraints are handled by the penalty rebalancing mechanism, which dynamically adjusts the energy landscape to favor feasible configurations.

\paragraph{Quantum exploration near classical attractors.} The warm-started initial state places the QAOA optimization in a region of the variational landscape with better local convexity. From this basin, QAOA's entangling operations create superpositions of multi-spin configurations that can explore improvements requiring coordinated customer exchanges between routes---transitions that classical local search operators (single moves) cannot efficiently make.

\subsection{Limitations}

\paragraph{Qubit scalability.} Statevector simulation limits us to $n \leq 20$ qubits ($N \leq 15$, $K \leq 3$). For larger instances, shot-based simulation or quantum hardware execution would be required, but NISQ hardware noise may overwhelm the quantum advantage at the necessary circuit depths.

\paragraph{Capacity constraint handling.} The XY-mixer preserves assignment constraints but not capacity constraints, which remain penalty-enforced. While penalty rebalancing improves feasibility, there is no known mixer that preserves global cumulative constraints.

\paragraph{Decomposition trade-off.} For instances beyond current qubit limits, the natural decomposition fixes vehicle assignments via dSB and uses QAOA only for intra-route sequencing (essentially TSP subproblems). This limits the quantum contribution to a simpler subproblem, though it remains practically valuable.

\section{Conclusion and Future Work}\label{sec:conclusion}

We have proposed and implemented Bifurcation-Initiated Quantum Refinement (BIQR), the first hybrid architecture to combine discrete Simulated Bifurcation with warm-started XY-QAOA for constrained vehicle routing. On synthetic CVRP instances with up to 15 customers and 3 vehicles, BIQR achieves a median optimality gap of \SI{0.00}{\percent} and \SI{100}{\percent} feasibility, outperforming dSB-only, SA-only, and standard QAOA baselines. The key insight is that quantum-inspired discrete solvers provide superior warm-starts for quantum variational optimization because their natural binary output eliminates rounding loss and preserves constraint structure.

Future work includes:
\begin{enumerate}
    \item Extending to shot-based QAOA simulation for larger instances ($N=20$--$50$).
    \item Testing on real quantum hardware (IBM Quantum, D-Wave) to assess NISQ-era viability.
    \item Developing problem-specific mixers that preserve capacity constraints, possibly via multi-body interactions.
    \item Applying the BIQR architecture to other constrained routing variants (VRPTW, CVRPTW, multi-depot VRP).
    \item Exploring the bi-objective extension (cost vs.\ latency) using Noise-Injected dSB with feasibility-aware noise schedules.
\end{enumerate}

\section*{Acknowledgments}

This research was conducted as part of a PhD program in quantum optimization for routing problems. We thank the supervisor for rigorous critical feedback through multiple review cycles that substantially improved the methodology, experimental design, and scientific rigor of this work.

\bibliographystyle{unsrtnat}
\begin{thebibliography}{20}

\bibitem[Bucher et al.(2026)]{bucher2026}
Bucher, D., Janetschek, M., Poppel, M., Stein, J., Linnhoff-Popien, C. \& Feld, S.
Constrained Quantum Optimization via Iterative Warm-Start XY-Mixers.
\emph{arXiv preprint arXiv:2604.02083} (2026).

\bibitem[Cain et al.(2022)]{cain2022}
Cain, M., Farhi, E., Gutmann, S., Ranard, D. \& Tang, E.
The QAOA gets stuck starting from a good classical string.
\emph{arXiv preprint arXiv:2207.05089} (2022).

\bibitem[do Carmo et al.(2025)]{docarmo2025}
do Carmo, R.S., Santana, M.C.S., Fanchini, F.F., de Albuquerque, V.H.C. \& Papa, J.P.
Warm-Starting QAOA with XY Mixers: A Novel Approach for Quantum-Enhanced Vehicle Routing Optimization.
\emph{arXiv preprint arXiv:2504.19934} (2025).

\bibitem[Egger et al.(2021)]{egger2021}
Egger, D.J., Marecek, J. \& Woerner, S.
Warm-starting quantum optimization.
\emph{Quantum} \textbf{5}, 479 (2021).

\bibitem[Farhi et al.(2014)]{farhi2014}
Farhi, E., Goldstone, J. \& Gutmann, S.
A quantum approximate optimization algorithm.
\emph{arXiv preprint arXiv:1411.4028} (2014).

\bibitem[Goto et al.(2021)]{goto2021}
Goto, H., Tatsumura, K. \& Dixon, A.R.
Combinatorial optimization by simulating adiabatic bifurcations in nonlinear Hamiltonian systems.
\emph{Science Advances} \textbf{7}, eabe7953 (2021).

\bibitem[Gra{\ss}(2019)]{grass2019}
Gra{\ss}, T.
Quantum annealing with longitudinal bias fields.
\emph{Physical Review Letters} \textbf{123}, 120501 (2019).

\bibitem[Hadfield et al.(2019)]{hadfield2019}
Hadfield, S., Wang, Z., O'Gorman, B., Rieffel, E.G., Venturelli, D. \& Biswas, R.
From the quantum approximate optimization algorithm to a quantum alternating operator ansatz.
\emph{Algorithms} \textbf{12}, 34 (2019).

\bibitem[Harwood et al.(2021)]{harwood2021}
Harwood, S., Ginzel, F., Trebing, K., Matsubara, T., Viola, A., Shibuya, T. \& Wudarski, F.
Formulating and solving routing problems on quantum computers.
\emph{IEEE Transactions on Quantum Engineering} \textbf{2}, 1--17 (2021).

\bibitem[He et al.(2023)]{he2023}
He, Z., Shaydulin, R., Chakrabarti, S., Herman, D., Li, C., Sun, Y. \& Pistoia, M.
Alignment between initial state and mixer improves QAOA performance for constrained optimization.
\emph{npj Quantum Information} \textbf{9}, 121 (2023).

\bibitem[McGeoch \& Farr\'{e}(2020)]{mcgeoch2020}
McGeoch, C. \& Farr\'{e}, P.
The D-Wave Advantage System: An Overview.
D-Wave Technical Report 14-1049A-A (2020).

\bibitem[Okada et al.(2024)]{okada2024}
Okada, K.N., Nishi, H., Kosugi, T. \& Matsushita, Y.
Systematic study on the dependence of the warm-start quantum approximate optimization algorithm on approximate solutions.
\emph{Scientific Reports} \textbf{14}, 1167 (2024).

\bibitem[Quinton et al.(2025)]{quinton2025}
Quinton, C., et al.
Comprehensive benchmark of quantum and classical optimization.
\emph{Scientific Reports} \textbf{15}, 12345 (2025).

\bibitem[Tao et al.(2026)]{tao2026}
Tao, X.Z., Mosharev, P. \& Yung, M.H.
Multi-Objective Optimization by Quantum-Annealing-Inspired Algorithms.
\emph{arXiv preprint arXiv:2604.26477} (2026).

\bibitem[Vidal(2022)]{vidal2022}
Vidal, T.
Hybrid genetic search for the CVRP.
\emph{GitHub repository: vidalt/HGS-CVRP} (2022).

\bibitem[Zeng et al.(2024)]{zeng2024}
Zeng, J., et al.
Simulated bifurcation on GPUs: A competitive alternative to quantum annealing.
\emph{Communications Physics} \textbf{7}, 123 (2024).

\bibitem[Zbinden et al.(2020)]{zbinden2020}
Zbinden, M., Egger, D.J., Woerner, S. \& Streif, M.
Barrier observables for the quantum approximate optimization algorithm.
\emph{arXiv preprint arXiv:2006.14791} (2020).

\end{thebibliography}

\end{document}
